\section{Perinatal pathways}

\subsection{Perinatal Regulatory Initialization}

PRI describes the transition from intrauterine regulation to respiration, feeding, temperature control, microbial acquisition, and social contact. Birth mode is an imperfect proxy for that multidimensional transition. Labor before caesarean delivery, surgical indication, gestational age, anesthesia, antibiotics, maternal illness, and postnatal care must be distinguished.

Human birth studies support large delivery-associated differences in vasopressin/copeptin. Wellmann et al. reported median arterial cord copeptin of 1,610 pmol/L after vaginal delivery versus 8 pmol/L after primary caesarean delivery in their cohort \cite{ref1}. This numerical contrast is specific to that cohort, assay, specimen, and time point. It establishes neither persistent hormonal deficiency nor a developmental effect of a specified magnitude.

A later prospective study explicitly found a birth-stress response for vasopressin but not oxytocin \cite{ref2}. The extension's broad grouping of all birth hormones must therefore be narrowed: AVP and OXT are distinct signals, and peripheral OXT is not a direct measurement of central affiliative circuitry. PRI predicts that timing and resolution of perturbation may matter; it does not imply that more stress is beneficial or that less stress is pathological.


\subsection{Maternal microbial inheritance}

Infant microbial assembly reflects multiple reservoirs and repeated exposures. Sinha et al.'s 2026 longitudinal study examined 714 mother–infant pairs and 4,526 samples, identifying the maternal gut as a major source of infant gut strains and associations with delivery and feeding \cite{ref3}. This supports a transmission ecology rather than a vaginal-canal-only model.

The causal variable proposed here is acquisition and persistence of particular organisms or functions, measured longitudinally. “Caesarean” and “vaginal” cannot substitute for those measurements. Antibiotics, diet, gestational maturity, household contacts, illness, and feeding can influence both microbial assembly and subsequent development. Stool non-detection is also not equivalent to biological absence at all mucosal sites.

Claims that a fixed percentage of modern humans has lost \emph{L. reuteri} are omitted because comparable populations, assay sensitivity, specimen sites, and strains have not been harmonized in this synthesis. A decline narrative is unnecessary for testing whether a particular strain or metabolic function affects a defined outcome.


\subsection{\emph{L. reuteri}, vagal signaling, oxytocin, and social reward}

In several mouse models, Sgritta et al. linked \emph{L. reuteri}-associated changes in social behavior to an intact vagus, oxytocin signaling, and social-interaction-related plasticity in ventral tegmental area circuitry \cite{ref4}. This is experimentally informative preclinical evidence. It is not proof that the same sequence explains human autism or that normal human development requires one commercially available strain.

Human findings remain narrower. The SB-121 crossover study enrolled 15 autistic participants aged 15–45 years, with safety as its primary objective \cite{ref5}. Adaptive-behavior signals were exploratory; the reported eye-tracking social-preference contrast had p = 0.1, and the plasma OXT contrast had p = 0.106. Neither constitutes a conventionally significant demonstration of the proposed human mediator. Multiple exploratory outcomes and formulation-specific effects limit interpretation.

A separate pediatric pilot trial reported improved social behavior without improvement in overall autism severity \cite{ref6}. This is compatible with dimensional specificity but does not establish developmental origin, vagal mediation, or \ensuremath{\alpha_2\delta} involvement. The key distinction is between a human intervention signal and verification of the full animal-derived pathway. A negative global autism endpoint must remain visible beside any positive social endpoint.


\subsection{Behavioral immune and disgust calibration}

The extension proposes two partially separable dimensions: affiliative/social-reward sensitivity and defensive pathogen-avoidance sensitivity. A schematic response model is Social approach = \ensuremath{\alpha} \ensuremath{\times} affiliative gain − \ensuremath{\beta} \ensuremath{\times} avoidance gain + contextual terms. The coefficients are not estimated and the dimensions need not be opposites; affiliation and pathogen avoidance can both be high.

The proposed bridge from early microbial/immune conditions to later germ aversion is E0. The Howard et al. meta-analysis distinguishes perceived infectability from germ aversion and reports different relationships with disgust and outgroup perceptions \cite{ref7}. It does not test PRI. Culture, threat learning, social experience, and measurement overlap provide competing explanations.

Disgust sensitivity and germ-aversion calibration are therefore a secondary horizon track. It must first predict age-appropriate contamination-avoidance measures before any broader social interpretation is attempted. Birth mode, microbiome status, diagnosis, or receptor state cannot be used to infer an individual's ideology, morality, or attitudes toward other people.
